\cotefile{157-4}
\scanpage{66}
\mspage{59}
% archive sheet number: 65
\uline{Théorème-scholie} : Soit $(C,\mathrm{TC},F,\mathrm{TF})$ une \unsure{quillénisation} uniforme gauche de $W$ sur $\mathrm{Diag}\subset\mathrm{Cat}$.

1) Soit $u:I'\unsure{-|\to}I$ une flèche de $\mathrm{Diag}$. On \deleted{\unsure{supposera que}} \unsure{$\mathrm{Diag}$ accordé|$\mathrm{Diag}$ associé} \unsure{à|:} $M$ de façon que pour $u$, flèche de $\mathrm{Diag}$, les $\varprojlim$ dans $M$ \deleted{\unsure{requises}} (de type $I'_{/i}$) requises pour assurer l'existence de
\[
u_*:M(I')\longrightarrow M(I)
\]
existent dans $M$. (C'est OK si $I^{\deleted{\unsure{'|*}}},I'$ finis, $M$ étant \unsure{supposé} stable par $\varprojlim$ finies de toute façon.) On sait (prop. 1 page 45) que $u_*$ induit
\[
u^{f}_{*\deleted{\unsure{f}}}:M(I')_f\longrightarrow M(I)_f
\tag*{(*)}
\]
On dit que $u$ est \deleted{\unsure{$W$-}}\uline{admissible} \inserted{(à gauche)} pour la quillénisation uniforme gauche donnée \inserted{de $W$ sur $\mathrm{Diag}$} si $u^f_*$ est compatible avec les localiseurs $W(I')_f$, $W(I)_f$ i.e. si $f':X'\unsure{-|\to}Y'$ dans $M(I')$, $X',Y'$ \deleted{$\unsure{\in}$} fibrants, $f'\in W(I')$ implique $u_*\unsure{(f)|(f')}\in W(I)$.

On dit que $Z$ dans $\mathrm{Diag}$ est admissible (pour la quillénisation uniforme gauche donnée de $W$ sur $\mathrm{Diag}$\ldots) si le foncteur $Z\unsure{-|\to}e$ est admissible. Ceci posé,

a) Pour que $u$ soit admissible, il suffit que les $I'_{/s}$ pour $s\in\mathrm{Ob}\,I$, le soient.
\margin{NB il y a une réciproque, quand $I$ est \deleted{$\unsure{I|s\backslash}$} ordonné, donc les \deleted{$\unsure{I|s\backslash}$} $I'_{/s}\unsure{-|\to}I'$ des immers. ouvertes.}

\scanpage{67}
\mspage{60}
% archive sheet number: 66
b) Tout ens. ordonné fini est admissible. Plus généralement (sauf erreur) tt ens. ordonné $I$ tel que l'espace topologique associé \unsure{à} $I^{\unsure{\circ|o}}$ soit \unsure{noeth.} (i.e. \deleted{toute} suite croiss. \unsure{d'ouverts} de $I^\circ$ est stationnaire, ou encore tte suite croissante de fermés de $I$ est stationnaire), et si de plus tte partie fermée irréductible de $I^\circ$ admet un pt générique, i.e. $I^\circ$ est (non seulement un espace noethérien, mais aussi) un espace top. \uline{sobre}. (Pour $I$, cela signifie que \deleted{\unsure{toute}} partie ouverte filtrante croissante de $I$ admet un plus grand élément\unsure{,|.}\deleted{\unsure{)}})

\deleted{\unsure{Ceci implique}}
c) En particulier, si $I',I$ \deleted{\unsure{font des} $\mathrm{Diag}$.} sont des ens. ordonnés finis, \unsure{toute flèche|les hom} $u:I'\unsure{-|\to}I$ est admissible.
\inserted{$=$, plus généralement, si $I^\circ,I'^\circ$ sont des \deleted{\unsure{espaces}} \inserted{ens. ordonnés} \deleted{\unsure{noeth. et sobres}} et $I'^\circ$ est un espace noeth. sobre.}

2) Supposons $u:I'\unsure{-|\to}I$ admissible, \deleted{\unsure{alors}} \deleted{\unsure{précisons que les} $\unsure{u_*|I}$} considérons
\diagram{
$u_{\unsure{*|\circ}}^W:$
\begin{tikzcd}[ampersand replacement=\&,column sep=huge,row sep=large]
M(I')_fW(I')_f^{-1} \arrow[r,dash] \arrow[d,"\unsure{q}"',"\approx"] \& \deleted{\unsure{M}}M(I)_fW(I)_f^{-1}\arrow[d,"\approx"] \\
\mathrm{HO}_W(I') \& \mathrm{HO}_W(I)
\end{tikzcd}
}
induit par $u_*$. Ce foncteur est adjoint à droite de
\[
u_W^*:\mathrm{HO}_W(I)\unsure{-|\longrightarrow}\mathrm{HO}_W(I').
\]

\scanpage{68}
\mspage{61}
% archive sheet number: 67
De plus, l'hom. canonique (pour $s\in\mathrm{Ob}\,I$, \inserted{$X'\in\mathrm{Ob}\,\mathrm{HO}_W(I')$})
\[
u_*^W(X')(s)\simeq p_{s*}\bigl(X'|_{\deleted{\unsure{s\backslash}}I'_{/s}}\bigr)
\overset{\text{def}}{=}
\underset{\deleted{\unsure{s\backslash}}I'_{/s}}{\underset{\longleftarrow}{\operatorname{\unsure{Holim|Rolim}}}}_W\bigl(X'|_{\deleted{\unsure{s\backslash}}I'_{/s}}\bigr)
\]
(où $p_s:\deleted{\unsure{s\backslash}}I'_{/s}\longrightarrow e$) est un iso, si on suppose que les $I'_{/s}$ sont admissibles \deleted{\unsure{(et $u$ admissible)}} \inserted{(\uline{strictement} admissible)}.

\uline{Corollaire} Sur le domaine formé des ens. ordonnés finis, \inserted{le prédérivateur} $\mathrm{HO}_W$ satisfait (en plus des axiomes $\mathrm{Der}\,1$, $\mathrm{Der}\,2$) aux axiomes $\mathrm{Der}\,4$ et $\mathrm{Der}\,5$. (existence des images directes $u_*$, adj. à droite des foncteurs $u^*$, et calcul \unsure{standard|standart} des fibres d'une image directe).

Démonstration du th.-scholie.

1°) a) est immédiat. Pour b) on procède par récurrence sur le cardinal $n$ de l'ens. ordonné fini $I$, les cas $n\in\{0,1\}$ étant triviaux. Supposons donc $n\geq2$ et le th. prouvé pour les $n'<n$. Je distingue deux cas : si $I$ admet un plus grand élém\textsuperscript{t} \deleted{\unsure{$b$}}\inserted{$a$}. Alors, désignant par $\alpha:e\unsure{-|\to}I$ le foncteur de valeur $a$, \inserted{par $p:I\unsure{-|\to}e$ la projection} on sait que
\[
p_*^{\deleted{\unsure{W|*}}}\simeq\alpha^{\unsure{*|\alpha}}\quad\text{i.e.}\quad\varprojlim_I X_i\simeq X_a,
\]
d'où le fait que sur $M(I)$ tout entier

\scanpage{69}
\mspage{62}
% archive sheet number: 68
(pas seulement sur $M(I)_f$) $p_*$ est compatible avec localiseurs $W(I)$, $W$.

Si $I$ n'a pas de plus grand élém\unsure{t|.}, alors l'ens. \inserted{$M$} de ses éléments maximaux est de card $\geq2$, on peut donc l'écrire $M=M_1^{\unsure{\circ|o|'}}\amalg M_2^\circ$, et on aura
\[
I=\unsure{I_1|\mathbb I_1}\cup I_2,\quad\text{où }I_1=I_{/M_1^\circ},\quad I_2=I_{/M_2^\circ}
\]
sont deux ouverts de $I$, \unsure{différents|distincts} de $I$. \inserted{donc de card $n_1,n_2<n$.} Considérons aussi $I_0=I_1\cap I_2$, c'est un ouvert de card. $<n$. Le diagramme
\diagram{
\begin{tikzcd}[ampersand replacement=\&,column sep=large,row sep=small]
\& I_1\arrow[dr,dotted] \& \\
I_0\arrow[ur]\arrow[dr] \& \& I=I_1\cup I_2 \\
\& I_2\arrow[ur,dotted] \&
\end{tikzcd}
}
correspond à une flèche
\diagram{
$u:I\longrightarrow\unsure{\Psi|\Upsilon}={}$
\begin{tikzcd}[ampersand replacement=\&,row sep=small]
\&1\\
0\arrow[ur]\arrow[dr]\&\\
\&2
\end{tikzcd},
}
avec
\[
I_1=I_{/1},\quad I_2=I_{/2},\quad I_0=I_{/0}.
\]
L'hyp. de récurrence et a) implique que $u$ est admissible. Comme un composé de foncteurs admissibles est \inserted{admissible} (tautologiquement, par $(vu)_*\simeq v_*u_*$), on est réduit finalement au

\scanpage{70}
\mspage{63}
% archive sheet number: 69
\uline{Lemme} La catégorie
\diagram{
$\unsure{\Psi|\Upsilon}={}$
\begin{tikzcd}[ampersand replacement=\&,row sep=small]
\&1\\
0\arrow[ur]\arrow[dr]\&\\
\&2
\end{tikzcd}
}
est admissible.

Ceci va résulter du

\uline{Cor.} Soit
\diagram{
$X_\bullet={}$
\begin{tikzcd}[ampersand replacement=\&,row sep=small]
\&X_1\arrow[dl,"\alpha_1"']\\
X_0\&\\
\&X_2\arrow[ul,"\alpha_2"]
\end{tikzcd}
}
dans $M(\Psi)$. Pour $X_\bullet$ soit fibrant, il f. et s. que

a) $X_0,X_1,X_2$ soient fibrants et

b) $\alpha_1,\alpha_2\in F$.
% margin note written slanted
\margin{\uline{NB} Il y a une variante évidente \unsure{caractérisant} les $X$ tels que $(X\unsure{-|\to}e)$ soit dans $\mathrm{TF}(\Psi)$, \unsure{comme ceux où|comme pour} $X_0,X_1,X_2$ sont \deleted{\unsure{f}}\unsure{fibrants} \uline{triviaux}, et $\alpha_1,\alpha_2\in\mathrm{TF}$.}

La nécessité de a) est \inserted{cas particulier} du cor. 2 page 48, celle de b) est cas part. du cor. 1 p 50.

Pour la suffisance, \inserted{\unsure{(dernier exercice préliminaire pour le th.)|(exercice préliminaire pour le th.)}} on \deleted{\unsure{doit|soit}} montre que dans le diagramme comm. à traits pleins
\diagram{
\begin{tikzcd}[ampersand replacement=\&,column sep=small,row sep=large]
A_{\unsure{1|0|\rho}}\arrow[rr]\arrow[dr,"\mathrm{TC}"']\arrow[dd] \& \& A_0\arrow[dr,"\mathrm{TC}"]\arrow[dd] \& \& A_2\arrow[ll]\arrow[dr,"\mathrm{TC}"]\arrow[dd] \& \\
\& B_1\arrow[rr]\arrow[dl,dashed,"h_1"'] \& \& B_0\arrow[dl,dashed,"h_0"'] \& \& B_2\arrow[ll]\arrow[dl,dashed,"h_2"] \\
X_1\arrow[rr,"\alpha_1"'] \& \& X_0 \& \& X_2\arrow[ll,"\alpha_2"] \&
\end{tikzcd},
}
avec trois flèches marquées $\mathrm{TC}$, on peut trouver $h_1,h_0,h_2$ laissant le diagr. comm. (trois triangles et deux carrés comm. rajoutés). Pour ceci, on commence par choisir $h_0$ de façon à rendre le triangle vertical médian comm. --- c'est possible car $X_0$ est fibrant, $A_0\unsure{-|\to}B_0$ dans $\mathrm{TC}$.

\scanpage{71}
\mspage{64}
% archive sheet number: 70
Ensuite, on voit séparément pour $h_1,h_2$ qu'on peut les trouver en \unsure{satisfaisant|astreignant} chacun aux deux contraintes de commutativité (1 triangle vertical et un carré penché). Pour $h_1$, p.ex., \unsure{l'existence} se voit sur le diagramme carré auxiliaire (traits pleins)
\diagram{
\begin{tikzcd}[ampersand replacement=\&,column sep=large,row sep=large]
A_1\arrow[r,"\mathrm{TC}","\deleted{\unsure{i}}"']\arrow[dd,dash] \& B_1\arrow[dd]\arrow[ddl,dashed,"h_1"']\arrow[dr,dashed,"\alpha_1^B"] \& \\
\&\&B_0\arrow[dl,dashed,"h_0"]\\
X_1\arrow[r,"F","\alpha_1^X"']\&X_0\&
\end{tikzcd}
}
\uline{NB} C'est la dém. de la prop. 3 p. 50

Le corollaire (partie \unsure{``il faut''}) étant vu, le \deleted{corollaire} lemme en résulte, car on sait que si on a un hom. $f$ de carrés cartésiens $X_\bullet$, $Y_\bullet$
\diagram{
\begin{tikzcd}[ampersand replacement=\&,column sep=large,row sep=large]
\&X_1\arrow[dl,"p_1^X"']\arrow[dd,"f_1"']\&\&X_3\arrow[ll]\arrow[dl]\arrow[dd,"f_3"]\\
X_0\arrow[dd,"f_0"']\&\&X_2\arrow[ll,"p_2^X"]\arrow[dd,"f_2"]\&\\
\&Y_1\arrow[dl,"p_1^Y"']\&\&Y_3\arrow[ll]\arrow[dl]\\
Y_0\&\&Y_2\arrow[ll,"p_2^Y"]\&
\end{tikzcd}
}
avec $p_{\unsure{1|2}}^X,p_2^X,p_1^Y,p_2^Y$ \unsure{fibrants|fibrantes} (il suffit $=$ $p_1^X$ et $p_1^Y$ \unsure{fibrants|fibrantes}), et $f_0,f_1,f_2\in W$ alors $f_3\in W$.

(\unsure{Ça|Ca} a été vu dans la théorie des produits fibrés $W$-homotopiques, et \deleted{\unsure{la}} vérification

\scanpage{72}
\mspage{65}
% archive sheet number: 71
directe est immédiate. De \deleted{fait que} \inserted{l'hyp \deleted{\unsure{inutile|utilisée}} pertinente,} que
\diagram{
\begin{tikzcd}[ampersand replacement=\&,column sep=small,row sep=small]
\&X_1\arrow[dl]\\
X_0\&\\
\&X_2\arrow[ul]
\end{tikzcd}
}
et
\diagram{
\begin{tikzcd}[ampersand replacement=\&,column sep=small,row sep=small]
\&Y_1\arrow[dl]\\
Y_0\&\\
\&Y_2\arrow[ul]
\end{tikzcd}
}
sont fibrants, on n'a besoin d'utiliser qu'une partie \unsure{dont|donc|seult}, viz. que $p_1^X$ et $p_1^Y$ sont fibrants, qui assure\deleted{\unsure{nt|t}} que ces deux carrés $X,Y$ sont $W$-homotopiquement cartésiens, ce qui suffit à impliquer dès que $f_0,f_1,f_2\in W\Rightarrow f_3\in W$.

Je vais essayer de généraliser ce raisonnement, pour traiter le cas d'un ens. ordonné $I$ qui peut être infini, mais tel que $I^\circ$ soit noethérien et sobre, donc toute \deleted{partie} suite croissante de parties fermées de $I$, et toute suite décroissante de parties ouvertes, est stationnaire, et de plus toute partie ouverte filtrante croissante a un plus grand élément.

On veut prouver que $I$ est admissible, et par ``récurrence noethérienne'' on est ramené à prouver ceci: $(\alpha)$ Si $I$ est
% vertical marginal emphasis line beside the condition
\margin{$\Big|$}
tel que pour tout ouvert $I'\subset I$, avec $I'\ne I$, \deleted{\unsure{on}} $I'$ soit admissible, alors $I$ aussi.

\unsure{Ceci|Cela} \deleted{\unsure{ceci}} \inserted{$(\alpha)$} admis, prouvons que $I\inserted{{}=\unsure{I_0|I_\circ}}$ est admissible.

\scanpage{73}
\mspage{66}
% archive sheet number: 72
S'il ne l'était, par $(\alpha)$ il existerait $I_1\subset I_0$, $I_1\ne I_0$ tel que $I_1$ non admissible. De $=$, il y en aurait $I_2\subset I_1$, $I_2\ne I_1$, non admissible etc. --- on pourrait construire une suite \inserted{str.} décroissante d'ouverts $I_n$ de $I$, contrairement au caractère noethérien de $I^\circ$.

Prouvons $(\alpha)$, et distinguons encore deux cas. a) $I^\circ$ est irréductible, donc admet un pt générique i.e. un plus petit élément, i.e. $I$ admet un plus grand élément. Ce cas a déjà été traité \unsure{tantôt} --- \deleted{toute} \unsure{tte} cat. avec objet final est admissible. b) $I^\circ$ est réductible, donc \deleted{$I^\circ$} réunion de deux fermés $I_1^\circ$, $I_2$ distincts de $I^\circ$, donc $I$ est réunion de deux ouverts $I_1,I_2$ distincts de $I$. Alors par l'hyp. \inserted{de} $(\alpha)$, $I_1,I_2$ \deleted{\unsure{et}} et $I_0=I_1\cap I_2$ sont admissibles. On termine alors comme tantôt, étant ramené à l'admissibilité de
\diagram{
\begin{tikzcd}[ampersand replacement=\&,column sep=small,row sep=small]
\&\&1\\
\unsure{\Psi|\Upsilon}=\&0\arrow[ur]\arrow[dr]\&\\
\&\&2
\end{tikzcd}
}
!

\scanpage{74}
\mspage{67}
% archive sheet number: 73
Prouvons la partie c) du th.-scholie, partie $1^\circ$. Cela résulte de a), une fois vu que pour $s\in\unsure{I|\mathbb{I}}$, $\deleted{\unsure{s\backslash}}I'_{/s}$ satisfait aux $=$s hypothèses que $I'$. Or $I$ \deleted{\unsure{étant|est}} ordonné, $\deleted{\unsure{s\backslash}}I_{/s}$ est une sous-catégorie ouverte de $I$, donc $I'_{/s}$ une sous-cat ouverte de $I'$, donc \inserted{$(I'_{/s})^\circ$} un sous-ens ordonné fermé de $I'^\circ$. Ce dernier étant noethérien et sobre, il en est de $=$ des \unsure{dits} fermés, d'où la conclusion.

Prouvons la partie $2^\circ)$ du th.-scholie. Considérons
\diagram{
\begin{tikzcd}[ampersand replacement=\&,column sep=large,row sep=large]
M(I')_c\arrow[r,hook]\&M(I')\&M(I')_f\arrow[l,hook]\arrow[d,"u_*^f"']\\
M(I)_c\arrow[r,hook]\arrow[u,"u_c^*"']\&M(I)\arrow[u,"u^*"']\&M(I)_f\arrow[l,hook]
\end{tikzcd}
}
,
où \deleted{$\unsure{u^*}$} \inserted{$u^*$ et a fortiori $u_c^*$} est compatible avec $W(I),W(I')$, et par hyp. $u_*$ est \unsure{comp.|comp} aux $W(I')_f$, $W(I)_f$. En passant aux cat. de fractions, on trouve

\scanpage{75}
\mspage{68}
% archive sheet number: 74
\diagram{
\begin{tikzcd}[ampersand replacement=\&,column sep=large,row sep=large]
M(I')_cW(I')_c^{-1}\arrow[r,dashed,"\simeq"]\&\deleted{\unsure{M(I')}}\mathrm{HO}_W(\unsure{I'|I})\&M(I')_f\arrow[l,dashed,"\simeq"']\arrow[d,"u_*^W"]\\
M(I)_cW(I)_c^{-1}\arrow[u,"u_W^{\unsure{*|\deleted{c}*}}"]\&\mathrm{HO}^{\deleted{\unsure{f|\circ}}}_W(I)\arrow[l,dashed,"\simeq"']\&M(I)_f\arrow[l,dashed,"\simeq"']
\end{tikzcd}
}
Et on veut montrer que $u_W^*$ et $u_*^W$ sont adjoints, de façon précise trouver un isom. bifonctoriel, pour \uline{$F$ dans $M(I)_c$} et \uline{$G'$ dans $M(I')_f$} \deleted{\unsure{iso}}
\[
\operatorname{Hom}_W(u^*F,\unsure{C'|G'})\xrightarrow{\sim}\operatorname{Hom}_W(F,u_*G')
\tag*{($\alpha$)}
\]
(où $\operatorname{Hom}_W$ désigne les $\operatorname{Hom}$ de $\mathrm{HO}_W(I')$ resp $\mathrm{HO}_W(I)$). Comme $F$ est dans $M(I)_c$, $u^*F$ est dans $M(I')_c$, et comme $G'$ est dans $M(I')_f$, $u_*(G')$ est dans $M(I)_f$. Il s'ensuit que les deux mbres de $(\alpha)$ sont des quotients des deux mbres dans la bijection canonique
\[
\operatorname{Hom}_{M(I')}(u^*F,G')\overset{\alpha}{\xrightarrow{\sim}}\operatorname{Hom}_{M(I)}(F,u_*G').
\tag*{($\alpha\alpha$)}
\]
Il faut simplement montrer que cette dernière transforme \unsure{la relation|les relations} d'équivalence pertinente $R'$ du premier membre,

\scanpage{76}
\mspage{69}
% archive sheet number: 75
en celle $R$ du second.

Je dis que l'on a de \unsure{ttes|toutes} façons
\[
R\leq\alpha(R'),
\]
\uline{sans} hypothèses d'admissibilité sur $u$, i.e. que la bijection inverse de $(\alpha\alpha)$ donne par passage aux quotients un hom.
\[
\operatorname{Hom}_W(u^*F,\unsure{G'|C'})\longleftarrow\operatorname{Hom}_W(F,u_*(G')),
\tag*{($\alpha$ bis)}
\]
forcément surjectif, et \inserted{de plus} fonctoriel en $F$, \deleted{$\unsure{G'}$} variables dans $\mathrm{HO}_W(I)$ et $\mathrm{HO}_W(I')$ respectivement). En effet, on définit la flèche $(\alpha\ \text{bis})$ comme composée
\diagram{
\begin{tikzcd}[ampersand replacement=\&,column sep=large,row sep=9em]
\operatorname{Hom}_W(F,u_*G')\arrow[r,"{\begin{gathered}\text{car }u^*\text{ est}\\\text{comp. avec }W(I),\\W(I')\text{ sur }M(I)\\\deleted{\text{tout}}\text{ entier, et}\\\text{on interprète les}\\\text{deux mbres comme des}\\\operatorname{Hom}\text{ dans }M(I)W(I)^{-1}\\\text{et }M(I')W(I')^{-1}\end{gathered}}"']\&\operatorname{Hom}_W(u^*F,u^*u_*G')\arrow[dr,"{\begin{gathered}\text{composition}\\\text{avec l'hom}\\\text{d'adjonction}\\u^*u_*G\to G\end{gathered}}"]\&\\
\&\&\operatorname{Hom}_W(u^*F,G')
\end{tikzcd}
}
Cela \deleted{signifie} implique donc que le foncteur \inserted{sur $\mathrm{HO}_W(I)$}
\[
F\longmapsto\operatorname{Hom}_{\mathrm{HO}_W(I')}(u^*F,G')
\]
est, sinon représentable, du moins quotient terme à terme du foncteur représenté par $u_*G'$. $^*$

% margin note written slanted
\margin{$^*$ Mais il n'est pas clair, sans hyp. suppl. sur la situation, que ce foncteur représentable puisse être choisi en dépendance fonctorielle de $G'$, envisagé dans $\mathrm{HO}_W(I')$.}

\scanpage{77}
\mspage{\unsure{70|78}}
% archive sheet number: 76
Mais supposons maintenant que $u$ soit admissible. Alors je vais définir la flèche en sens initial (p. 68), \deleted{de la façon suivante}: Je choisis une flèche
\[
\deleted{(\alpha)\ \operatorname{Hom}_W(u^*F\unsure{,|)}}
\]
\[
i'':u^*F\longrightarrow F'
\]
avec $i''\in\mathrm{TC}\subset C\cap W$, $F'\in\operatorname{Ob}M(I')_f$, (donc en fait \deleted{\unsure{de ce fait}} $F'\in\operatorname{Ob}M(I')_{cf}$ car $u^*F$ est cofibrant, \unsure{comme on sait|comme quoi})

\deleted{ou ce qui revient au $=$ \unsure{par}} et par adjonction la flèche $i'$
\[
i':F\unsure{-|\longrightarrow}u_*F'.
\]
Je considère le diagramme (comm. !)
\diagram{
\begin{tikzcd}[ampersand replacement=\&,column sep=huge,row sep=large]
\operatorname{Hom}(F',\unsure{G'|C'})\arrow[r,"\alpha"]\arrow[d,"\lambda:"',"f'\mapsto f'\circ i''"]\&\operatorname{Hom}(u_*F',u_*G')\arrow[d,"\mu:f\mapsto f\circ i'"']\\
\operatorname{Hom}(u^*F,G')\arrow[r,"\mathrm{adj.}"']\&\operatorname{Hom}(F,u_*G')
\end{tikzcd}
}
On voit que \deleted{$\unsure{F'}$} l'on a encore $F'$ cofibrant (car $u^*F$ l'est et $\unsure{i'|i''}\in\mathrm{TC}\subset C$) donc $\operatorname{Hom}_W(F',G')$ est un quotient de $\operatorname{Hom}(F',G')$, et $\lambda$ induit un isomorphisme
\[
\operatorname{Hom}_W(F',G')\xrightarrow{\sim}\operatorname{Hom}_W(u^*F,G')
\]
(car $i''\in\mathrm{TC}\deleted{\unsure{)|f}}\subset W$). Comme $F',G'\in\operatorname{Ob}M(I')_f$, et que $u_*^f$ est compatible avec $W(I')_f$, $W(I)_f$, il \unsure{s'ensuit} que la flèche horizontale $\alpha$

\scanpage{78}
\mspage{71}
% archive sheet number: 77
est comp. avec les relations d'équivalences définies par
\[
\operatorname{Hom}(F',\unsure{G'|C'})\unsure{-|\longrightarrow}\operatorname{Hom}_{W(I')}(F',G'),
\]
resp
\[
\operatorname{Hom}(u_*F',u_*G')\longrightarrow\operatorname{Hom}_{W(I)}(u_*F',u_*G').
\]
Et il en est de $=$ de la flèche $\mu:f\mapsto f\circ i'$, trivialement. Donc la flèche d'adjonction $\operatorname{Hom}(u^*F,G')\longrightarrow\operatorname{Hom}(F,u_*G')$ est compatible avec les relations $R,R'$, ce qui achève de prouver que $(\unsure{\alpha\alpha|\alpha})$ page 68 induit une bijection $(\alpha)$ p. 68.

\deleted{\unsure{les}} Ceci montre donc $1^\circ)$ que pour un objet fixé $G'$ dans $\mathrm{HO}_W(I')$ (qu'on peut tjrs supposer, à isom près, provenir de $M(I')_f$) le foncteur
\[
F\longmapsto\operatorname{Hom}_{\mathrm{HO}_W(I')}(u_W^*F,G')
\]
est représentable --- i.e. le foncteur $u_W^*$ admet un adjoint à D., et $2^\circ)$ que cet adjoint à gauche se calcule comme $u_*(G')$
\margin{\unsure{Quel allonge\uline{ment}}}
($G'$ est pris fibrant) --- \deleted{comme} vu que l'hom $(\alpha\ \text{bis})$ page 69 est \uline{fonctoriel en $F$}. (A posteriori il s'ensuit qu'il est aussi fonctoriel en $G'$\ldots)

\scanpage{79}
\mspage{72}
% archive sheet number: 78
Le dernier énoncé du $2^\circ$ du th.-scholie résulte du calcul de $u_{\unsure{*|\circ}}^W$, compte tenu que $\deleted{\unsure{i'|I'}\longleftarrow\unsure{s|*}}\ \deleted{\unsure{s\backslash}}I'_{/s}\unsure{-|\longrightarrow}I'$ est un iso local, donc le \inserted{foncteur} image inverse \deleted{\unsure{d'un|d'une} \unsure{est|et}} envoie
\[
M(I')_f\unsure{-|\longrightarrow}M(I'_{/s})_f,
\]
donc si $\unsure{G'|C'}$ est dans $M(I')_f$,
\[
u_*(G')(s)\simeq p_{s*}(G'|I'_{/s})
\]
est un calcul de $p_{s*}^W(G'|I'_{/s})$.

Cela achève la démonstration du th.-scholie sur les images directes.

\uline{Commentaire.} 1) Ce théorème me paraît loin d'être satisfaisant, c'est un pis-aller. J'avais espéré, avec les hypothèses du \S\ 3, qu'on aurait l'existence du $u_*^W$ pour \uline{toute} flèche $u:I'\to I$ dans Cat. Je doute à présent que ce soit le cas, mais n'ai pas encore \unsure{trouvé|essayé} de faire un contre-exemple. Les cas les plus importants \inserted{pour la question d'admissibilité d'une catégorie $I$} me paraissent les suivants:

a) Catégorie \inserted{$I$} discrète infinie (disons dénombrable). La question revient ici à celle-ci: soient $f_i:X_i\to Y_i$ dans $W$, avec $X_i,Y_i$ fibrants donc dans $F$. Alors $f=\prod f_i:\prod X_i\to\prod Y_i$ est-il dans $W$?

\scanpage{80}
\mspage{73}
% archive sheet number: 79
b) $I$ catégorie
\diagram{
\begin{tikzcd}[ampersand replacement=\&,column sep=tiny]
0\&\unsure{-|\to}\&1\&\unsure{-|\to}\&2\&\unsure{-|\to}\&3\arrow[r]\&\ldots
\end{tikzcd}
}
i.e. la catégorie ordonnée associée à l'ens ordonné $\mathbb{N}$, pour la relation d'ordre ordinaire $m\leq n$. Donc $M(I)$ est formé des systèmes \uline{projectifs} dans $M$, indexés par $\mathbb{N}$. \deleted{\unsure{On}} On prouve facilement \inserted{(en procédant comme pour le cas du lemme p. 63)} que $M(I)_f$ est formé des syst. projectifs
\diagram{
\begin{tikzcd}[ampersand replacement=\&,column sep=small]
X_0\&X_1\arrow[l]\&X_2\arrow[l]\&\ldots\arrow[l]\&X_n\&X_{n+1}\arrow[l,"\unsure{p_n|\rho_n}"']
\end{tikzcd}
}
tels que les $X_i$ soient dans $M_f$, et les $p_n$ dans $F$. La question est donc: si $X_\bullet,Y_\bullet$ sont fibrants et $f_\bullet:X_\bullet\to Y_\bullet$ est tel que $f_i\in W_f$ pour $i\in\mathbb{N}$, alors a-t-on $\varprojlim f_i:X_\infty\to Y_\infty$ dans $W$?

Si \deleted{\unsure{la}} la réponse est affirmative \inserted{pour $\mathbb{N}$}, (pour un $W$ et une quillénisation \unsure{donné|donnés|donnée}), elle l'est aussi pour \deleted{toute} catégorie discrète dénombrable, comme on le voit \unsure{aussitôt}. Mais je ne vois pas comment déduire, $=$ \unsure{ça des axiomes|des axiomes|de ces axiomes}, disons pour un triplet de Quillen clos (closed model category) tel que $M$ soit de plus stable par $\varinjlim$ filtrantes, que celles-ci soient exactes, qu'il y ait

\scanpage{81}
\mspage{74}
% archive sheet number: 80
un petit ensemble de gén. de $M$ (car $M$ est accessible), que $W$ soit stable par $\varinjlim$ filtrantes. Je devrais demander à Thomason \unsure{s'il|s'îl} \unsure{connait|connaît} un contreexemple.

c) Les deux exemples précédents étaient des catégories ordonnées. Comme exemple typique de cat. non ordonnée, où je ne vois pas ce qui se passe, je signale $I=B_{\unsure{G|C}}$, $G$ un groupe (\unsure{donc} un seul objet, avec un groupe d'autom. $G$). Aussi le cas de $\unsure{\mathbb{\Delta}|\Delta}[1]=\{\mathbb{\Delta}^0,\mathbb{\Delta}^1\}$ (sous-cat. pleine de $\mathbb{\Delta}$) n'est pas clair pour moi.

2) On aimerait pouvoir traiter de façon duale l'existence des $u_!^W$, à défaut de mieux. Dans l'hypothèse de départ du \S\ 4 sur $(M,W)$, cela consiste donc à changer \unsure{des $+$ et des $-$|des $+$ en $-$}. Cette fois, on devra supposer l'existence d'une quillénisation donc \uline{exacte à droite} pour $W$. Ce qui est plus exigeant, semble-t-il, qu'une telle structure exacte à gauche, dans le cas envisagé), avec
\[
\mathrm{TF}=F\cap C,
\]
